\documentclass[10pt,a4paper]{amsart}
\usepackage[margin=19mm]{geometry}
\usepackage{amsmath,amssymb,mathtools}
\usepackage[scr=rsfs,cal=euler]{mathalfa}
\usepackage{xfrac}
\usepackage[unicode,hidelinks,bookmarks=false]{hyperref}
\newcommand{\ie}{i.e.\ }
\newcommand{\cf}{cf.\ }
\newcommand{\resp}{respectively\ }
\newcommand{\Subsection}{\subsection}
\providecommand{\nobreakdash}{\nobreak\hbox{-}}
\setlength{\emergencystretch}{2em}
\multlinegap0pt
\usepackage{bm}
\usepackage{bbm}
\usepackage{dsfont}
\usepackage{xspace}
\usepackage{cancel}
\usepackage{mathtools}
\usepackage{amsthm}
\makeatletter
\@ifundefined{theorem}{\newtheorem{theorem}{Theorem}}{}
\@ifundefined{lemma}{\newtheorem{lemma}[theorem]{Lemma}}{}
\@ifundefined{proposition}{\newtheorem{proposition}[theorem]{Proposition}}{}
\makeatother
\newcommand{\rng}{{\rm rng} \hspace{0.03cm}}
\title[Bounded ranges of cardinal functions]{Bounded ranges of cardinal functions}
\author{Jacek Marchwicki, B{\l}a\.zej \.Zmija}
\date{Source: arXiv:2508.13016v1 (CC BY 4.0)}
\begin{document}
\maketitle

\noindent\emph{Source and licence.} Bounded ranges of cardinal functions. Version arXiv:2508.13016v1 (CC BY 4.0), distributed under the Creative Commons Attribution 4.0 International licence, as recorded in the arXiv licence metadata for this exact version and shown on the abstract page https://arxiv.org/abs/2508.13016v1. Full rights notice: licenses/arxiv-2508.13016-CC-BY-4.0.txt.\par\smallskip

\noindent\textit{Editorial scope.} Counted unit: the Proposition whose source label is \texttt{PropSetsI\_1} in the article (source0.tex lines 1049--1051, source label \texttt{PropSetsI\_1}), delivered together with the article's own complete original proof of it (source0.tex lines 1052--1111, one \texttt{proof} environment of 60 lines). No mathematics is written, repaired or re-derived: statement and proof are the article's own text line for line. Required context retained uncounted: lematdodwojkowych (lines 395--400); mniejnizszesc (lines 1018--1029). Every definition, display and label of those lines is retained unchanged. Omitted from this package and not redistributed: 1--394, 401--1017, 1030--1048, 1112--1793. Nothing inside a retained excerpt is omitted or altered: every retained body is delivered exactly as the article's lines are written, so no omission receipt of any kind accompanies this package. Citations: the retained excerpts carry no citation command at all, so no bibliography entry of the article is transcribed and no reference is redistributed. Reference adaptation: none was needed and none was made; every reference inside the delivered unit resolves inside it, so no unresolved reference or citation is produced. Printed slips kept literal: no printed word, symbol, hypothesis or number of the counted statement or of its proof was altered. Layout and class: the article's own class is \texttt{amsart} with packages including amssymb, inputenc, picture, color, tikz, pifont, amsthm, enumitem; the wrapper is the lane's pinned \texttt{amsart} editorial wrapper at 10pt on A4, which restates the theorem-like environments the retained text needs and adds the article's own macro definitions that the retained text needs (\texttt{\textbackslash rng}), with extra wrapper packages (bm, bbm, dsfont, xspace, cancel, mathtools). The delivered numbering is the unit's own. Assets: no retained span contains a figure environment, a graphics-inclusion command, a PDF-page inclusion, a table or a TikZ picture; no image file is copied into this package, the package declares no asset dependency and no image file is redistributed with it. Licence: version arXiv:2508.13016v1 of this article is distributed under the Creative Commons Attribution 4.0 International licence, as recorded in the arXiv licence metadata for this exact version and shown on the abstract page https://arxiv.org/abs/2508.13016v1. A full rights notice is delivered at licenses/arxiv-2508.13016-CC-BY-4.0.txt. Characters: every retained excerpt is pure ASCII, so the delivered engine, XeLaTeX with its default Latin Modern text font, reports no missing character.
\begin{lemma}\label{lematdodwojkowych}
Let $x_n=\frac{1}{2^n}$ and $f$ be the cardinal function for $(x_n)$. Then 
\begin{displaymath}
f(x) = \left\{ \begin{array}{ll} 2, & \textrm{if $x\in  \mathcal{D}\setminus\{0,1\}$,}\\ 1, & \textrm{if $x\in \big((0,1)\setminus \mathcal{D}\big)\cup\{0,1\}$.} \end{array} \right.
\end{displaymath}
\end{lemma}

\begin{lemma}\label{mniejnizszesc}
Let $\mathbf{x}=(x_n)$ be a sequence with cardinal function $f$. Let $z$ and $w$ be positive real numbers. Assume that one of the following holds:
\begin{enumerate}
\item $z+w>\sum_{n=1}^{\infty}x_n$,
\item $z=w= \frac{1}{2}\sum_{n=1}^{\infty}x_n$.
\end{enumerate}
Let $h$ be the cardinal function of the sequence $z\uplus (w\uplus \textbf{x}) = (z,w,x_{1},x_{2},\ldots )$. Then 
\begin{align*}
\max (h)\leq 3\max (f).
\end{align*}
In particular, if $\rng (f)=\{1,2\}$ then $\max (h)\leq 6$.
\end{lemma} 

\begin{proposition}\label{PropSetsI_1}
$\{1,2,3\}$, $\{1,2,3,4\}$, $\{1,2,3,4,5\}$, $\{1,2,3,4,6\}$, $\{1,2,3,4,5,6\}\in\mathcal{I}_{1}$.
\end{proposition}

\begin{proof}
For every set $M$ from the statement, we will find a sequence $\mathbf{x}$ with the cardinal function $h$ such that $\rng (h)=M$. At first, let $\mathbf{b}= \left(\frac{1}{2^{n}}\right)$ and $f$ be its cardinal function. In the constructions, we will consider sequences of the form $\mathbf{x} = z\uplus (w\uplus \mathbf{b})$ for several choices of $z$ and $w$. Note that $\rng (f)=\{1,2\}$ by Lemma \ref{lematdodwojkowych}, so $\rng (h)\leq 6$ by Lemma \ref{mniejnizszesc}.

Now let us proceed case by case.
\begin{enumerate}
	\item For $\{1,2,3\}$: let $z>w$ be non-dyadic numbers such that $z\in (1,1+w)$, $w\in(\frac{1}{2},1)$, and $z-w$ is non-dyadic. Then clearly $\mathcal{A}(\mathbf{x}) = [0,1+z+w]$ and $1,2\in \rng (h)$. Moreover, observe that for every $x$:
	\begin{align*}
	h(x) = f(x) + f(x-z) + f(x-w) + f(x-z-w),
	\end{align*}		
	and at most one of the numbers $x$, $x-z$, $x-w$, $x-z-w$ is dyadic. Moreover, the relations between sizes of $z$ and $w$ imply that always at most two of these these numbers lie in the interval $[0,1]$. Indeed,
	\begin{itemize}
	\item if $x\in [0,1]$ then $x-z<0$ and $x-z-w<0$,
	\item if $x\in (1,1+w]$ then $x-z-w<0$ and $x>1$,
	\item if $x\in [1+w,1+z+w]$ then $x>1$ and $x-w>1$.
	\end{itemize}	
	 Therefore, Lemma \ref{lematdodwojkowych} yields $h(x)\leq 3$. Observe, that also $h(x)=3$ for every dyadic $x\in (0,1)$. Hence, $\rng (h)=\{1,2,3\}$.
	\item For $\{1,2,3,4\}$: here we use a similar choice, as in the previous construction. The only change is that now we require $z-w$ to be dyadic. Then at most two of the numbers $x$, $x-z$, $x-w$, $x-z-w$ can be dyadic, so $h(x)\leq 4$. Moreover, $h(x)=3$ for every dyadic $x\in (w,1)$ and $h(x)=4$ for example for every $x=z+\frac{1}{2^{m}}$, where $m$ is any number such that $z+\frac{1}{2^{m}} < 1+w$. Hence, $\{1,2,3,4\}\in\mathcal{I}_{1}$.
	\item For $\{1,2,3,4,5\}$: let us take $z=w\in (\frac{1}{2},1)$ be any non-dyadic number. Then $\mathcal{A}(\mathbf{x})=[0,1+2z]$, $1,2\in\rng (h)$, and
	\begin{align*}
	h(x)=f(x)+2\cdot f(x-z)+f(x-2z).
	\end{align*}
	At most on of the values $x$, $x-z$, $x-2z$, so
	\begin{align*}
	h(x)\leq \max\{2+2\cdot 1, 1 +2\cdot 2\}=5.
	\end{align*}
	It remains to show, that $h(x)$ can be equal to $3$, $4$ and $5$.
	\begin{itemize}
		\item $h(x)=3$ for $x=z$.
		\item $h(x)=4$ for $x=1-\frac{1}{2^{m}}$, where $m$ is any positive integer such that $z<1-\frac{1}{2^{m}}$,
		\item $h(x)=5$ for $x=z+\frac{1}{2^{m}}$, where $m$ is such that $z+\frac{1}{2^{m}} < 1$.
	\end{itemize}
	\item For $\{1,2,3,4,6\}$: let $z=w=\frac{1}{2}$. Clearly $\mathcal{A}(\mathbf{x})=[0,2]$. Moreover, for every $x$:
	\begin{align*}
		h(x) = f(x)+2\cdot f\left(x-\frac{1}{2}\right)+f(x-1).
	\end{align*}
	Using the above equality, we can write down a formula for $h(x)$. For $x\in[0,1]$ we have:
	\begin{align*}
	h(x) = \left\{
	\begin{array}{ll}
		1, & \textrm{ if } x=0, \\
		2, & \textrm{ if } x\in (0,\frac{1}{2})\cap \mathcal{D},\\
		3, & \textrm{ if } x\in [0,1]\setminus \mathcal{D},\\
		4, & \textrm{ if } x=\frac{1}{2},\\
		6, & \textrm{ if } x\in (\frac{1}{2},1]\cap\mathcal{D}.
	\end{array}	
	\right.
	\end{align*}
	If $x\in (1,\frac{3}{2})$ then $h(x)=2f\left(x-\frac{1}{2}\right) +f(x-1)$. We have that either both, $x-\frac{1}{2}$ and $x-1$, are dyadic, or non of them. Therefore, in this case $h(x)\in \{3,6\}$. Then $h\left(\frac{3}{2}\right)= 2f(1) + f\left(\frac{1}{2}\right) = 4$. Finally, if $x\in (\frac{3}{2},2]$ then $h(x) = f(x-1)\in\{1,2\}$. Therefore, $\rng (h) =\{1,2,3,4,6\}$.
	\item For $\{1,2,3,4,5,6\}$: we take $z=w\in\ (\frac{1}{2},1)$ to be any dyadic number. Lemma \ref{mniejnizszesc} yields $\max (h)\leq 6$. Moreover, using the same reasoning, as in the proof of the case $\{1,2,3,4,5\}$, we get
	\begin{itemize}
		\item $h(0)=1$,
		\item $h(x)=2$ for every dyadic $x\in (0,z)$,
		\item $h(x)=3$ for every non-dyadic $x\in (z,1)$, 
		\item $h(z)=4$,
		\item $h(1)=5$,
		\item $h(x)=6$ for $x=z+\frac{1}{2^{m}}$, where $m$ is a positive integer such that $z+\frac{1}{2^{m}} < 1$.
	\end{itemize}
	Thus $\rng (h) = \{1,2,3,4,5,6\}$.
\end{enumerate}
\end{proof}

\end{document}
